%% THIS has been moved from the other paper

\CatchFileBetweenTags{\CabibboAngleVal}{constants.tex}{CabibboAngleVal}
\CatchFileBetweenTags{\CabibboAngleExperimentalValue}{constants.tex}{CabibboAngleExperimentalValue}
\CatchFileBetweenTags{\CabibboAngleEq}{constants.tex}{CabibboAngleEq}
\CatchFileBetweenTags{\CabibboAngleSigma}{constants.tex}{CabibboAngleSigma}

\subsection{The Leakage Partition (Cabibbo Angle \texorpdfstring{$\theta_C$}{thetaC})}
\label{subsec:generation_coupling}
\label{subsec:cabibbo}

\subsubsection{The Architectural Requirement}
Under \textbf{Rule 1: Capacity Partitioning}, a persistent system must allocate fractional bandwidth to handle internal state transitions without disrupting its primary operations. The third \textbf{Protocol} allocation of the Capacity Partition is the \textbf{Generation Coupling}: the strictly bounded geometric leakage rate between the resonant tiers (generations) of the lattice. In the Standard Model, this manifests as the Cabibbo Angle ($\theta_C \approx 13.0^\circ$, or the primary CKM matrix element $|V_{us}|$), parameterizing quark mixing between the first and second generations.

We derive the Cabibbo Angle geometrically as this specific fractional bandwidth allocation. While the Weak Mixing Angle partitions the macroscopic spacetime forces, the Cabibbo Angle represents the strictly bounded \textbf{Geometric Leakage} of state transitions within the internal generational structure.

To derive this leakage probability, we must define the exact geometric ratio of the inter-generational interface to the available active channel capacity using the substrate invariants:

\begin{enumerate}
    \item \textbf{The Interface ($\pi$):} Generational mixing is an internal state transition, occurring within the node's symmetry space prior to macroscopic spatial propagation. It is therefore bounded by the \textbf{unscaled} geometry of the fundamental interaction loop (Wilson Loop), not the macroscopic boundary ($\pi\Delta$).

    On the discrete $D_4$ lattice, the Wilson Loop is a polygon. However, generational mixing is a \textbf{resonance transition} between distinct generational states of the same topological defect. The defect's internal structure is defined by the continuous field approximation (hydrodynamic limit, where wavelength $\gg$ lattice spacing). In this limit, the loop perimeter becomes the continuum value $\pi d$ for diameter $d=1$.

    The value $\pi$ is forced: the mixing angle must be transcendental to avoid rational commensurability with the lattice periodicity. If $\sin\theta_C$ were rational, the generational states would be resonantly locked to the lattice, preventing independent evolution. $\pi$ is the unique transcendental arising from isotropic circle geometry in the continuous limit.

    \item \textbf{The Active Flavor Width ($\nu - \chi$):} The total chiral capacity is $\nu=16$. However, the topological boundary $\chi=2$ serves as the ``Identity Lock'' that maintains the particle's structural stability against the vacuum. These 2 channels are \textbf{not available} for mixing: if they participated, the defect would lose its identity, becoming flavor violation rather than generational mixing.

    The free bandwidth is therefore strictly $\nu - \chi = 14$. Using $\nu = 16$ would predict $\sin\theta_C = \pi/16 \approx 0.196$, failing by $4\sigma$. The boundary lock is geometrically mandatory; its subtraction is not a fitting choice but the conservation of topological identity.
\end{enumerate}

\subsubsection{The Leakage Ratio}
The mixing angle is defined by the strict geometric ratio of the Interface to the Active Flavor Width. This determines the maximum angle at which a signal can ``slip'' from one generation to the next without exceeding the available channel capacity and breaking the topological lock. 

Physically, this ratio represents the quantum tunneling amplitude between generational resonances: while the boundary ($\chi$) protects the particle's core identity, the open channels ($14$) permit a precise, geometrically defined leakage rate between adjacent generation tiers.

\begin{equation}
\sin \theta_C = \frac{\text{Interface}}{\text{Flavor Width}} = \frac{\pi}{\nu - \chi}
\end{equation}

No other ratio is possible: $\pi$ is the unique continuum-limit circumference, and $\nu - \chi = 14$ is the unique free capacity after identity lock. The aspect ratio $\pi/14 \approx 0.224$ is the maximal geometric leakage that preserves both (i) transcendental mixing (incommensurate with lattice periodicity) and (ii) topological identity (boundary lock intact).

\textbf{Numerical Calculation:}
\begin{equation}
\sin \theta_C = \frac{\pi}{14} \approx \mathbf{\CabibboAngleVal}
\end{equation}
\begin{equation}
\theta_C = \arcsin\left(\frac{\pi}{14}\right) \approx 12.97^\circ
\end{equation}

\textit{Scope Defense:} The leakage ratio uses the bare, unscaled boundary ($\pi$) because generational mixing is strictly an internal state transition. It occurs within the node's internal symmetry space prior to macroscopic spatial propagation, bounded by internal hardware capacity ($14$) and unscaled geometry. The macroscopic boundary $\pi\Delta$ governs spatial propagation (System 5, Section \ref{sec:resonant_interface}); here we address internal mixing (System 4, Capacity Partition).

\subsubsection{Validation: Precision Flavor Physics}
We compare this geometric derivation to the world average extracted from kaon and hyperon decays.

\vspace{0.5em}
\begin{tabular}{ll}
\textbf{Prediction:} & \CabibboAngleVal \\
\textbf{Experiment:} & \CabibboAngleExperimentalValue \\
\textbf{Accuracy:}   & $\CabibboAngleSigma \sigma$ \\
\end{tabular}
\vspace{0.5em}

\noindent \textbf{Physical Interpretation:} The Cabibbo Angle is exactly the geometric aspect ratio ($\pi/14$) of a chiral channel constrained by a topological boundary, representing the maximum permissible information leakage rate of the lattice.

\subsubsection{Geometric Consequence: The Structural Unity of Gauge and Flavor}

The derived Cabibbo mixing ($\sin \theta_C \approx 0.224$) is nearly identical to the Weak Mixing Angle ($\sin^2 \theta_W \approx 0.223$). In the Standard Model, this is an unexplained empirical coincidence. Here it is a \textbf{structural necessity}: both are partition coefficients of the same finite bandwidth budget.

This near-equality, often noted as an unexplained empirical coincidence in the Standard Model, is here revealed as a structural consequence: both quantities are partition coefficients of the exact same finite bandwidth budget. The geometry that partitions the forces ($\Delta/N_{sys}$) is structurally coupled to the geometry that mixes the masses ($\pi/(\nu-\chi)$). Both are manifestations of the exact same underlying lattice capacity, rigidly allocating finite bandwidth between Temporal Coherence (the forces, $\Delta/N_{sys}$) and Generational Leakage (the matter, $\pi/(\nu-\chi)$). $E_8$-Persistence theory provides the structural mechanism uniting gauge and flavor sectors: the same lattice capacity $\nu=16$ rigidly allocates bandwidth between Temporal Coherence ($\Delta/N_{sys}$) and Generational Leakage ($\pi/(\nu-\chi)$). The near-equality is not accidental; it is the signature of a single, exhausted budget.